\chapter{Trabalhos Relacionados}
\label{cap:trabalhos-relacionados}

Este capítulo reúne estudos sobre a segmentação de exsudatos e a classificação da retinopatia diabética em imagens de fundoscopia. A busca foi feita no IEEE Xplore, no PubMed e no Google Acadêmico. A \textit{string} de busca usada foi: (\mbox{\textit{``diabetic retinopathy''}} OR \mbox{\textit{``fundus image''}}) AND (\mbox{\textit{``exudate segmentation''}} OR \mbox{\textit{``hard exudates''}} OR \mbox{\textit{``soft exudates''}} OR \mbox{\textit{``lesion segmentation''}} OR \mbox{\textit{``grading''}} OR \mbox{\textit{``classification''}}) AND (\mbox{\textit{``deep learning''}} OR \mbox{\textit{``U-Net''}} OR \mbox{\textit{``Transformer''}}). Os resultados foram limitados a artigos em inglês publicados entre 2020 e 2026 em periódicos com revisão por pares. A restrição a periódicos deixa de fora conferências importantes da área, como a MICCAI, e foi adotada para manter trabalhos com a descrição completa do método e da avaliação. Foram incluídos os estudos que segmentam exsudatos ou outras lesões da RD, ou que classificam a RD, com aprendizado profundo, usando bases públicas, como IDRiD, DDR, APTOS e Messidor, e descrevendo o protocolo de avaliação, com prioridade para os que tratam de exsudatos moles e duros, comparam arquiteturas ou propõem uma solução CAD. Ao final, foram selecionados onze trabalhos, a maior parte publicada em periódicos da IEEE, como o \textit{IEEE Access} e o \textit{IEEE Transactions on Medical Imaging}, e os demais em revistas da área, como \textit{Diagnostics}, \textit{Electronics}, \textit{Disease Markers} e \textit{International Journal of Imaging Systems and Technology}. Por ser uma revisão direcionada, e não sistemática, a quantidade de resultados de cada base não foi contabilizada, e a seleção não tem a pretensão de ser exaustiva. A parte de classificação reúne apenas três trabalhos, e comparações entre CNNs e \textit{Transformers} em imagens médicas publicadas fora de periódicos, como a de \citeonline{matsoukas2021cnnvit}, ficaram fora da seleção por esse critério e aparecem apenas na Fundamentação Teórica.

A apresentação segue duas linhas. A Seção~\ref{sec:rel-segmentacao} descreve as pesquisas de segmentação de lesões da retina, desde variantes da U-Net até modelos com atenção, redes leves e \textit{Transformers}, e a Seção~\ref{sec:rel-classificacao} descreve os estudos que classificam a RD a partir da imagem inteira. Para cada trabalho, são apresentados o método, as bases utilizadas, o principal resultado e a limitação mais relevante para esta proposta. Ao final, a Seção~\ref{sec:rel-sintese} compara os trabalhos entre si e com esta proposta.

\section{Segmentação de exsudatos e lesões retinianas}
\label{sec:rel-segmentacao}

\citeonline{zong2020unet} propõem, no \textit{IEEE Access}, uma U-Net com módulos \textit{Inception} e conexões residuais para a segmentação de exsudatos duros. Antes de chegar à rede, os autores recortam o fundo preto, trabalham com o canal verde e realçam o contraste; depois, usam o algoritmo de \textit{superpixels} SLIC para gerar recortes de $256 \times 256$ pixels, como forma de lidar com o tamanho reduzido do conjunto e com o desbalanceamento entre classes. Na inferência, as predições dos recortes vizinhos são combinadas para recompor a máscara da imagem completa. A perda escolhida é a \textit{focal loss}, para que os pixels de fundo, muito mais numerosos, não dominem o treinamento. Avaliado no IDRiD, com a partição oficial de 54 imagens de treino e 27 de teste, o método alcança sensibilidade de 96,38\%, especificidade de 97,14\% e acurácia de 97,95\%.

As ablações apresentadas pelos autores indicam que o recorte guiado pelo SLIC funciona melhor que o recorte aleatório e que a combinação de \textit{Inception} com conexões residuais rende mais do que cada componente sozinho. A principal limitação está no protocolo, que considera apenas exsudatos duros, uma única base, nenhum teste em outro domínio e métricas pixel a pixel que, com tanto fundo na imagem, tendem a ficar altas e dificultam a comparação com trabalhos que usam a AUPR do desafio do IDRiD \cite{porwal2020idrid}.

\citeonline{ma2024hard} também segmentam exsudatos duros com uma U-Net, mas em um cenário com poucos dados, o DiaRetDB1, com apenas 89 imagens originais, das quais as máscaras binárias são geradas a partir das anotações em XML da própria base, isolando somente as regiões de exsudato duro. A rede segue o esquema clássico de codificador e decodificador, com entrada de $256 \times 256$ pixels, e o estudo utiliza aumento de dados espacial (rotações, deslocamentos, cisalhamento e espelhamentos), sem alterar a cor dos pixels, para não descaracterizar as lesões. O treinamento usa entropia cruzada binária e otimizador Adam, com \textit{early stopping} para conter o sobreajuste, algo especialmente importante quando quase todo o volume de treino vem de cópias transformadas das mesmas imagens. Com validação cruzada de seis partições, os autores reportam F1 de 0,9754. Esse valor deve ser interpretado com cautela, pois a base é pequena, apenas os exsudatos duros são considerados e os próprios autores reconhecem que faltam validação em conjuntos maiores e comparação com outros métodos para confirmar o desempenho.

\citeonline{kou2020erunet} apresentam, também no \textit{IEEE Access}, a ERU-Net, voltada a microaneurismas e exsudatos. A arquitetura parte de uma U-Net residual e organiza três caminhos de reconstrução que começam em profundidades diferentes, buscando recuperar detalhes que se perdem quando a lesão é pequena ou pouco contrastada; o treino é feito em recortes de $48 \times 48$ pixels sobre o canal verde realçado. Um ponto forte é a avaliação em três bases, E-ophtha, IDRiD e DDR, com validação cruzada de cinco partições no E-ophtha e no IDRiD. Para exsudatos, a AUC-ROC chega a 0,9866 no IDRiD e a 0,9609 no DDR, acima da U-Net e da ResU-Net no mesmo protocolo e com menos parâmetros que esta última. Os autores ainda testam a arquitetura na segmentação do disco óptico, indicando que o desenho multicaminho pode ser aplicado a outros tipos de estrutura. O estudo também reforça a escolha do IDRiD e do DDR para a etapa de segmentação deste trabalho. Contudo, a AUC-ROC, em um problema no qual quase todos os pixels pertencem ao fundo, tende a indicar um desempenho otimista. Também se observa que cada base é avaliada separadamente, sem treinar em uma e testar em outra, e os exsudatos moles não são considerados.

\citeonline{yin2023dbunet} apresentam a DB-U-Net, uma rede de dois caminhos voltada a lesões pequenas. O primeiro caminho é generalista, e o segundo prevê ao mesmo tempo a lesão e a sua borda, com dois decodificadores ligados ao mesmo codificador; um módulo de fusão com ASPP junta tudo no final, e a \textit{Dice loss} é usada contra o desbalanceamento. Os testes passam por IDRiD, DDR e E-ophtha, com AUPR como métrica principal. No IDRiD, o método alcança AUPR de 0,7297 para exsudatos moles, o melhor valor nessa classe no comparativo apresentado pelos autores, e de 0,8878 para exsudatos duros. Os experimentos de ablação indicam que tanto o decodificador de contorno quanto a fusão dos dois ramos contribuem para esse ganho. Para os autores, a supervisão das bordas favorece a segmentação de lesões pequenas e de forma irregular. Os próprios autores, no entanto, reconhecem que o método ainda oscila diante de variações de iluminação, resolução e tamanho das lesões. Por outro lado, manter duas U-Nets completas em paralelo aumenta o custo de treino e de inferência. Isso dificulta o uso em cenários com pouca infraestrutura, e as bases continuam sendo avaliadas uma a uma, sem um teste explícito de generalização entre domínios.

\citeonline{wan2021eadnet} propõem, na revista \textit{Disease Markers}, a EAD-Net, com codificador residual, um módulo de atenção dual (de posição e de canal) e decodificador no estilo U-Net, treinada com uma combinação de entropia cruzada binária e Dice. O mesmo modelo segmenta as quatro lesões clássicas da RD, incluindo exsudatos duros e moles, e por isso constitui um precedente de segmentação multilesão, sobretudo porque várias das equipes mais bem colocadas no desafio do IDRiD treinavam uma rede diferente para cada lesão \cite{porwal2020idrid}. No pré-processamento, as bordas pretas são removidas, as imagens são normalizadas para $1024 \times 1024$ pixels e o treino ocorre em $512 \times 512$, com espelhamentos, mudanças de escala, translações e rotações como aumento de dados. A avaliação usa o e\_ophtha\_EX, subconjunto de exsudatos da base E-ophtha, com 82 imagens, o IDRiD e uma base local de 262 fundoscopias anotadas por oftalmologistas. No e\_ophtha\_EX, a rede alcança sensibilidade de 92,77\%, e no IDRiD os autores reportam AUPR comparável à dos melhores colocados do desafio. A comparação, porém, exige cuidado. No e\_ophtha\_EX as métricas são calculadas pelo casamento entre candidatos detectados e lesões anotadas, e não pixel a pixel, e a base local não é pública, o que impede reproduzir essa parte dos resultados.

\citeonline{ali2023lezioseg} propõem o LezioSeg na revista \textit{Electronics}, um modelo leve no qual os autores usam uma MobileNet pré-treinada no ImageNet como codificador, um módulo ASPP para contexto em várias escalas, conexões de salto com portas e um bloco de atenção multiescala no decodificador, além de aumento de dados por transformações afins. O modelo resultante tem cerca de 10,7 milhões de parâmetros, menos que muitas das alternativas comparadas. No IDRiD, com 54 imagens de treino e 27 de teste, o modelo alcança AUPR de 86\% para exsudatos duros e de 81\% para exsudatos moles, no conjunto de comparação do próprio trabalho; no E-ophtha, chega a 63\% para exsudatos duros, e o DDR é usado para testar a generalização, sem retreinamento, o que fornece uma estimativa mais realista do desempenho fora do domínio de origem. O uso de arquiteturas leves é próximo do que se propõe neste trabalho, que inclui MobileNet e EfficientNet entre as redes avaliadas. Os autores aplicam o teste \textit{t} de Student, mas apenas para comparar a versão final com a própria rede de base. A limitação está no protocolo, pois todos os resultados vêm de uma única divisão fixa, sem validação cruzada. Com apenas 14 imagens de teste contendo exsudatos moles, não é possível saber o quanto esses números variariam com outra partição.

\citeonline{huang2022rtnet} apresentam, no \textit{IEEE Transactions on Medical Imaging}, o RTNet, um modelo com \textit{Transformers} projetado para segmentar ao mesmo tempo microaneurismas, hemorragias e exsudatos duros e moles. A proposta parte de uma observação clínica: as lesões não aparecem isoladas e costumam ter relação espacial com os vasos. Por isso, a rede combina uma U-Net com codificador DenseNet-161 a blocos \textit{Transformer} que modelam as relações entre lesões e entre lesões e vasos. Como não existe uma base pública com anotações conjuntas de vasos e lesões da RD, as máscaras de vasos usadas no treino são pseudo-rótulos gerados por um modelo treinado no DRIVE e no STARE. A avaliação é feita no IDRiD e no DDR, e no IDRiD o RTNet alcança AUPR de 0,9024 para exsudatos duros e de 0,7502 para exsudatos moles, os melhores valores nessas duas classes no comparativo apresentado pelos autores. Os autores ainda fazem um teste de generalização, treinando no DDR e testando no IDRiD.

Segundo os próprios autores, a baixa qualidade das pseudo-máscaras de vasos restringe o desempenho máximo da rede. O trabalho é uma referência para a segmentação conjunta de exsudatos duros e moles, mas o custo do codificador DenseNet-161 junto com os blocos de atenção dificulta o uso em ambientes com poucos recursos, e o desempenho é reportado em uma única partição por base.

\citeonline{zheng2024cascade} publicaram no \textit{International Journal of Imaging Systems and Technology} um modelo em dois estágios para as quatro lesões clássicas da RD (exsudatos duros e moles, hemorragias e microaneurismas), segmentadas de uma só vez. Cada estágio é uma U-Net modificada com blocos \textit{Transformer}: o primeiro gera uma máscara inicial e o segundo a corrige, combinando convoluções, que captam detalhes locais, com atenção, que capta o contexto da imagem. Os testes usam o IDRiD, o DDR e o DiaRetDB1, e os autores relatam desempenho competitivo em relação a outros métodos, com destaque para exsudatos moles e microaneurismas. Assim como a EAD-Net e o RTNet, esse trabalho combina U-Net e atenção. No desenho em dois estágios, porém, erros da primeira etapa podem se propagar para a segunda, e o custo computacional aumenta, já que são duas redes em sequência.

\section{Classificação da retinopatia diabética}
\label{sec:rel-classificacao}

\citeonline{naveen2025effnet} apresentam o EffNet-SVM no \textit{IEEE Access}, usando uma EfficientNetV2-Small pré-treinada no ImageNet para extrair características da fundoscopia e um SVM com \textit{kernel} RBF no lugar da camada \textit{softmax}. São comparadas três configurações: a EfficientNet com \textit{softmax}, um SVM alimentado com descritores tradicionais e o híbrido proposto. O estudo usa o APTOS 2019, mas agrupa os cinco graus em um problema binário (1.805 imagens sem RD e 1.857 com RD), com as imagens redimensionadas para $224 \times 224$ pixels. Nesse cenário, o modelo alcança acurácia de 97,26\%, à frente da própria EfficientNet com \textit{softmax} e de outras CNNs, como MobileNetV2, ResNet50 e DenseNet-121. Mapas Grad-CAM são usados para indicar as regiões que mais influenciam a decisão, embora a interpretabilidade obtida seja apenas parcial, já que o SVM continua pouco transparente. A proposta se relaciona com a ideia de CAD e com o uso de arquiteturas eficientes, mas o ponto mais relevante é a queda no teste externo: ao avaliar o modelo no HRF, com apenas 30 imagens (15 de olhos saudáveis e 15 com RD), os autores relatam acurácia de 76,76\%. Esse valor não corresponde a um número inteiro de acertos em 30 imagens e deve ser lido com cautela, mas a queda mostra o quanto o resultado no APTOS depende do domínio de origem, além de a tarefa binária deixar de lado a gradação clínica.

\citeonline{yang2024tmil} propõem, no \textit{IEEE Access}, o TMIL, uma forma de usar o ViT sem descartar a alta resolução das câmeras de fundo de olho. Em vez de reduzir a imagem inteira para $224 \times 224$ pixels, ela é redimensionada para $896 \times 896$ e dividida em 16 partes menores; cada parte passa pelo ViT e um bloco de atenção global reúne essa informação antes de decidir o grau da RD. No APTOS, em cinco classes, a versão final do modelo alcança acurácia de 85,6\%, kappa de 90,7\% e sensibilidade de 73,7\% (o resumo do artigo informa kappa de 91,5\%, valor obtido pelas versões intermediárias); no Messidor-1, em uma tarefa binária de RD referenciável, chega a 93,1\% de acurácia, reduzindo ainda o tempo de inferência em cerca de 62\% em relação ao ViT alimentado com a imagem inteira. Esse ganho vem do encurtamento da sequência de entrada, já que cada parte gera 196 \textit{tokens}, contra 3.136 quando a imagem de $896 \times 896$ é passada inteira ao ViT. O trabalho mostra que reduzir muito a resolução da imagem pode prejudicar a classificação com \textit{Transformers}. A limitação está na avaliação, feita em uma única divisão do APTOS, com o treino repetido cinco vezes, mas sem validação cruzada, e uma sensibilidade de 73,7\% que ainda deixa escapar parte dos casos doentes.

O trabalho mais próximo desta proposta é o de \citeonline{boumeshed2026modular}, publicado na revista \textit{Diagnostics}. Os autores montam uma solução modular com duas partes. A primeira segmenta as quatro lesões no IDRiD, com uma U-Net com codificador ResNet-34 para cada lesão, e a segunda classifica o grau da RD de 0 a 4 com uma EfficientNet-B0 treinada no APTOS 2019. O classificador depois é testado nas 1.744 imagens graduáveis do Messidor-2 sem nenhum ajuste, e a queda de desempenho é acentuada. O QWK cai de 0,873 no teste do APTOS para 0,455 no Messidor-2, e a sensibilidade para RD referenciável cai de 92,7\% para 35,0\%. A avaliação é mais cuidadosa que a da maioria dos trabalhos citados aqui, com intervalos de confiança por \textit{bootstrap}, quatro sementes de treino e uma checagem de imagens duplicadas no APTOS. Os autores também testam acrescentar ao classificador características extraídas das lesões segmentadas, mas isso não melhora a gradação de forma significativa. Na segmentação, o Dice nos recortes com lesão fica entre 0,51 e 0,68, enquanto o Dice por imagem, que inclui os recortes sem lesão, fica entre 0,21 e 0,61. Em relação a esta proposta, a limitação é que cada tarefa usa uma única arquitetura, então o estudo não mostra se uma rede leve ou baseada em \textit{Transformer} perderia menos desempenho fora do domínio de origem.

Para facilitar a comparação, a Tabela~\ref{tab:comparativo-relacionados} resume, para cada um dos onze trabalhos descritos neste capítulo, as bases utilizadas, a tarefa, a arquitetura, o tipo de validação, a presença de validação externa, de teste estatístico e de solução CAD, além da métrica principal relatada. A última linha descreve o que é proposto neste trabalho, e não resultados já obtidos.

{%
	\SingleSpacing%
	\setlength{\tabcolsep}{2.5pt}%
	\renewcommand{\arraystretch}{1.15}%
	\setlength{\LTcapwidth}{\textwidth}%
	\captionsetup{justification=raggedright,singlelinecheck=false}%
	\begin{longtable}{@{}>{\scriptsize\raggedright\arraybackslash}p{1.8cm}
		>{\scriptsize\raggedright\arraybackslash}p{1.8cm}
		>{\scriptsize\raggedright\arraybackslash}p{1.3cm}
		>{\scriptsize\raggedright\arraybackslash}p{2.5cm}
		>{\scriptsize\raggedright\arraybackslash}p{1.7cm}
		>{\scriptsize\centering\arraybackslash}p{0.8cm}
		>{\scriptsize\centering\arraybackslash}p{1.3cm}
		>{\scriptsize\raggedright\arraybackslash}p{2.2cm}
		>{\scriptsize\centering\arraybackslash}p{0.7cm}@{}}
		\caption{\label{tab:comparativo-relacionados} Comparação entre os trabalhos relacionados e este trabalho}\\
		\toprule
		Trabalho & Base(s) & Tarefa & Arquitetura & Validação & Val.\ ext. & Teste estat. & Métrica principal & CAD \\
		\midrule
		\endfirsthead
		\multicolumn{9}{@{}l}{\footnotesize\textit{Continuação da Tabela~\ref{tab:comparativo-relacionados}}}\\
		\toprule
		Trabalho & Base(s) & Tarefa & Arquitetura & Validação & Val.\ ext. & Teste estat. & Métrica principal & CAD \\
		\midrule
		\endhead
		\midrule
		\multicolumn{9}{r@{}}{\footnotesize\textit{Continua na próxima página}}\\
		\endfoot
		\bottomrule
		\multicolumn{9}{@{}p{15.3cm}@{}}{\ABNTEXfontereduzida Fonte: elaborada pelo autor.}\\
		\multicolumn{9}{@{}p{15.3cm}@{}}{\ABNTEXfontereduzida Nota: \checkmark~atende; $\triangle$~atende parcialmente (na coluna Val.\ ext., de validação externa, mais de uma base avaliada, mas sem treinar em uma e testar em outra; na coluna CAD, integração apenas discutida ou feita só como fusão de características, sem ferramenta de apoio; na coluna Teste estat., teste ou intervalo de confiança sem comparação entre arquiteturas diferentes); $\times$~não atende ou não relatado; N/I: protocolo de validação não identificado na leitura realizada. EX: exsudatos duros; SE: exsudatos moles; HE: hemorragias; MA: microaneurismas.}\\
		\endlastfoot
			\citeonline{zong2020unet} & IDRiD & EX & U-Net com \textit{Inception} e resíduo & \textit{Hold-out} (54/27) & $\times$ & $\times$ & Sens.\ 96,38\%; Acur.\ 97,95\% & $\times$ \\
			\citeonline{ma2024hard} & DiaRetDB1 & EX & U-Net modificada & 6-\textit{fold} & $\times$ & $\times$ & F1 0,975 & $\times$ \\
			\citeonline{kou2020erunet} & E-ophtha, IDRiD, DDR & EX, MA & ERU-Net (U-Net residual multicaminho) & \textit{Hold-out}; 5-\textit{fold} & $\triangle$ & $\times$ & AUC-ROC (EX) 0,987 IDRiD; 0,961 DDR & $\times$ \\
			\citeonline{yin2023dbunet} & IDRiD, DDR, E-ophtha & EX, SE, HE, MA & DB-U-Net (U-Net + U2Net + ASPP) & \textit{Hold-out} (54/27 no IDRiD) & $\triangle$ & $\times$ & AUPR IDRiD: EX 0,888; SE 0,730 & $\times$ \\
			\citeonline{wan2021eadnet} & e\_ophtha\_EX, IDRiD, base local (262) & EX, SE, HE, MA & EAD-Net (residual + atenção dual) & \textit{Hold-out} (54/27 no IDRiD) & $\triangle$ & $\times$ & Sens.\ 92,77\% (e\_ophtha\_EX) & $\times$ \\
			\citeonline{ali2023lezioseg} & IDRiD, E-ophtha, DDR & EX, SE, HE, MA & LezioSeg (MobileNet + ASPP) & \textit{Hold-out} (54/27) & \checkmark & $\triangle$ & AUPR IDRiD: EX 0,86; SE 0,81 & $\times$ \\
			\citeonline{huang2022rtnet} & IDRiD, DDR & EX, SE, HE, MA & RTNet (U-Net/DenseNet-161 + \textit{Transformer}) & \textit{Hold-out} & \checkmark & $\times$ & AUPR IDRiD: EX 0,902; SE 0,750 & $\times$ \\
			\citeonline{zheng2024cascade} & IDRiD, DDR, DiaRetDB1 & EX, SE, HE, MA & U-Nets em cascata com \textit{Transformer} & N/I & $\triangle$ & $\times$ & PR-AUC, Dice, IoU (destaque em SE e MA) & $\times$ \\
			\citeonline{naveen2025effnet} & APTOS, HRF & RD binária & EfficientNetV2-S + SVM & N/I & \checkmark & $\times$ & Acur.\ 97,26\% (APTOS); 76,76\% (HRF) & $\triangle$ \\
			\citeonline{yang2024tmil} & APTOS, Messidor-1 & Graus 0--4; binária & ViT + MIL (TMIL) & \textit{Hold-out} (10\% teste; 5 repetições) & $\triangle$ & $\times$ & Acur.\ 85,6\%; kappa 90,7\% (APTOS) & $\times$ \\
			\citeonline{boumeshed2026modular} & IDRiD; APTOS, Messidor-2 & EX, SE, HE, MA; graus 0--4 & U-Net (ResNet-34) por lesão; EfficientNet-B0 & \textit{Hold-out}; 4 sementes & \checkmark & $\triangle$ & Dice 0,51--0,68 nos recortes com lesão e 0,21--0,61 por imagem (seg.); QWK 0,873 (APTOS), 0,455 (Messidor-2) & $\triangle$ \\
		\midrule
			\textbf{Este trabalho (proposto)} & IDRiD, DDR; APTOS, EyePACS, Messidor-2 & EX, SE; graus 0--4 & U-Net única com 6 codificadores: VGG16, ResNet50, MobileNetV3-Large, EfficientNet-B0, ViT-B/16, Swin-T & Busca única + 5-\textit{fold} $\times$3 (seg.) / $\times$2 agrupado por paciente (clf.) & \checkmark & Wilcoxon e \textit{t} corrigido + Holm; IC \textit{bootstrap}; DeLong & AUPR, Dice (seg.); QWK, AUC referenciável (clf.); parâmetros, MACs e latência & \checkmark \\
	\end{longtable}%
}

\section{Síntese comparativa e lacuna abordada}
\label{sec:rel-sintese}

De modo geral, os trabalhos analisados se concentram em uma única tarefa ou arquitetura. Dos nove trabalhos que fazem segmentação, seis tratam exsudatos duros e moles juntos \cite{huang2022rtnet,yin2023dbunet,wan2021eadnet,ali2023lezioseg,zheng2024cascade,boumeshed2026modular}, e os outros três ficam nos exsudatos duros, com ou sem microaneurismas. Apenas dois dos onze estudos usam validação cruzada \cite{kou2020erunet,ma2024hard}. Só quatro testam o modelo em uma base diferente daquela usada no treino \cite{huang2022rtnet,ali2023lezioseg,naveen2025effnet,boumeshed2026modular}, e os demais avaliam cada conjunto separadamente ou ficam em uma única base. A análise estatística também é pouco frequente. \citeonline{ali2023lezioseg} usam o teste \textit{t} para comparar o LezioSeg com a sua rede de base, e \citeonline{boumeshed2026modular} trazem intervalos de confiança por \textit{bootstrap}, mas nenhum dos trabalhos compara arquiteturas diferentes com teste de hipótese e correção para várias comparações. As métricas também mudam muito de um trabalho para outro. Parte usa AUPR, parte usa AUC-ROC e parte usa métricas pixel a pixel, o que dificulta comparações diretas.

Quanto às arquiteturas, modelos leves aparecem em três trabalhos \cite{ali2023lezioseg,naveen2025effnet,boumeshed2026modular} e \textit{Transformers} em três \cite{huang2022rtnet,zheng2024cascade,yang2024tmil}, mas, entre os onze trabalhos selecionados, nenhum compara redes tradicionais, leves e baseadas em \textit{Transformers} sob o mesmo protocolo. O que mais se aproxima disso é o comparativo entre CNNs feito por \citeonline{naveen2025effnet}, ainda em uma tarefa binária. A gradação em cinco classes aparece em dois estudos \cite{yang2024tmil,boumeshed2026modular}. Só \citeonline{boumeshed2026modular} juntam segmentação de lesões e classificação na mesma solução, com uma arquitetura em cada tarefa e uma rede separada para cada lesão, e \citeonline{naveen2025effnet} apenas discutem essa possibilidade.

Diante disso, este trabalho propõe segmentar exsudatos moles e duros, classificar a retinopatia diabética em graus 0 a 4 e comparar seis modelos representativos das famílias tradicionais, leves e baseadas em \textit{Transformers} sob o mesmo protocolo, com os seis codificadores ligados a uma única U-Net, que segmenta os exsudatos duros e moles ao mesmo tempo, em vez de uma rede para cada lesão. A proposta inclui validação cruzada agrupada por paciente quando a base permite, validação externa no DDR e no Messidor-2, testes de Wilcoxon e \textit{t} corrigido com ajuste de Holm, intervalos \textit{bootstrap} e a medida do custo de cada rede, além de reunir os melhores modelos em uma solução de apoio ao diagnóstico que apresente a máscara de exsudatos, o grau previsto e um alerta de edema macular na mesma interface.
